Este capítulo cierra la memoria con un balance del trabajo realizado. Se revisan los objetivos planteados en la introducción, se recogen las lecciones técnicas y metodológicas que ha dejado el desarrollo, y se identifican las líneas de continuación que el proyecto deja abiertas.

\section{Objetivos alcanzados}

El objetivo principal de este proyecto era desarrollar una aplicación que permitiese a personas con dificultades del habla construir frases seleccionando pictogramas y convertirlas en voz natural. El trabajo realizado cubre los siete objetivos específicos de la introducción en los términos siguientes.

El capítulo de antecedentes recoge el análisis del estado del arte (objetivo 1): revisa los sistemas SAAC existentes, las técnicas de generación de lenguaje natural aplicadas a CAA y los trabajos previos en la intersección entre pictogramas y NLG. Esta revisión fundamentó las decisiones de diseño de las tres técnicas base y permitió identificar los patrones de hibridación explorados en la cuarta fase.

El vocabulario seleccionado para el objetivo 2 consta de 232 pictogramas ARASAAC organizados en nueve categorías cromáticas, acompañados del \textit{golden standard} de 164 casos de prueba con múltiples referencias humanas (sección~\ref{sec:eval-golden}), que cubre ocho categorías semánticas y tres niveles de complejidad. El corpus funciona como referencia común para la evaluación de los seis motores.

La implementación y evaluación comparativa de tres enfoques de IA (objetivo 3) constituye el núcleo del trabajo. El proyecto incluye tres motores de generación (gramática formal con Lark, embeddings semánticos con FastText y LLM con GPT-4o) junto con tres hibridaciones que los combinan: post-edición LLM sobre gramática (H1), post-edición LLM sobre embeddings (H2) y sobregeneración multimotor con ranking por perplejidad (H3). Los seis motores alcanzan una tasa de éxito del 100\% y una cobertura de pictogramas entre el 96\% y el 99\%. La evaluación comparativa, con once métricas automáticas, tests de significancia estadística y análisis por complejidad, caracteriza las fortalezas y compromisos de cada enfoque.

La aplicación móvil integra el motor TTS nativo de cada plataforma para la síntesis de voz (objetivo 4), lo que permite la lectura en voz alta de las frases generadas.

Dixi cubre el objetivo 5 (interfaz accesible): una aplicación desarrollada con React Native y Expo que presenta los pictogramas organizados por categorías cromáticas y permite construir frases mediante selección directa, priorizando la simplicidad de interacción para usuarios con necesidades de comunicación aumentativa.

La validación con hablantes nativos de español se articula en dos encuestas (objetivo 6). La primera construye el \textit{golden standard} multi-referencia contra el que operan las métricas automáticas: 57 participantes inscritos, 45 completaron los 20 ítems de su bloque y aportaron 948 respuestas que se consolidaron en 615 referencias únicas. La segunda mide la preferencia humana entre las tres técnicas base mediante elección forzada: 54 participantes, 48 completos, 1275 respuestas sobre 44 ítems. El LLM recibió el 44\% de las elecciones, con diferencias pareadas frente a gramática y embeddings que resisten la corrección de Bonferroni a $p < 10^{-17}$. La validación con usuarios finales de CAA y con profesionales del ámbito queda fuera del alcance del TFG como línea de continuación. La documentación técnica (objetivo 7) se recoge en los capítulos de desarrollo y en la presente memoria.

\section{Lecciones aprendidas}

El desarrollo de seis motores de generación sobre el mismo dominio acotado (frases cortas en español a partir de 232 pictogramas) ha producido aprendizajes que los resultados numéricos por sí solos no reflejan.

El liderazgo del LLM en la evaluación, tanto automática como humana, ha sido más claro de lo que el diseño del TFG anticipaba. El motor basado en GPT-4o domina las seis métricas con referencia humana (Exact Match del 63\%, BLEU-4 de 0.78, chrF de 0.85, BERTScore de 0.967) y recoge el 44\% de las elecciones en la encuesta de naturalidad, el doble que cualquier motor base, con diferencias pareadas que resisten la corrección de Bonferroni. La correlación entre chrF, BERTScore y Exact Match del LLM y la preferencia humana en los ítems solapados entre ambas encuestas ($\rho \approx 0.37\text{--}0.38$, $p < 0.05$) valida empíricamente el uso de estas métricas como indicadores de calidad percibida.

La gramática formal y los embeddings semánticos se equiparan en la práctica. Producen la misma salida en 163 de los 164 casos del \textit{golden standard} (99.4\%), sus métricas coinciden en la tercera cifra decimal y la encuesta de naturalidad no los distingue estadísticamente. El componente de embeddings, diseñado para resolver predicción de preposiciones y \textit{ranking} semántico, no altera las decisiones de la gramática en el vocabulario acotado del proyecto. Los embeddings pueden diferenciarse en escenarios con vocabulario más amplio donde las reglas léxicas no cubran todos los casos, una hipótesis que queda como línea abierta.

Los híbridos cumplen sus objetivos de diseño en las dimensiones específicas para las que fueron pensados pero no en coincidencia con la referencia humana. Los híbridos 1 y 2 reducen la perplejidad mediana a 70 (frente a 119 de las bases y 123 del LLM), lo que indica que la post-edición por LLM produce texto más natural según los modelos de lenguaje preentrenados. Hybrid3 alcanza la máxima cobertura (0.99) y fidelidad ajustada (0.986). En las métricas con referencia humana, los tres híbridos mejoran sobre las bases pero no alcanzan al LLM puro: la post-edición sobre una salida determinista añade valor, pero menos que dejar al LLM generar desde cero.

Hybrid3 constituye un resultado negativo instructivo. Su diseño (sobregeneración con tres motores y \textit{ranking} por perplejidad) es conceptualmente sólido y está respaldado por la literatura. En la práctica, selecciona la salida de la gramática en 120 de 164 casos (73.2\%), porque sus tres candidatas tienden a converger en la misma formulación y el ranking por perplejidad no tiene granularidad suficiente para discriminar entre ellas. La latencia de 4.3 segundos, la más alta de todos los motores, no se justifica ante una ganancia marginal en cobertura sobre Hybrid1. Este resultado confirma que la técnica adecuada depende del dominio: \textit{overgenerate-and-rank} funciona cuando las candidatas son diversas y la función de \textit{ranking} puede discriminar entre ellas; en un dominio con frases cortas y vocabulario restringido, ninguna de las dos condiciones se cumple.

La convergencia entre motores ha sido un hallazgo transversal revelador. En 113 de los 164 casos (68.9\%) el LLM genera de forma independiente exactamente la misma oración que la gramática, a pesar de ser un modelo generativo sin acceso a las reglas del CFG. El dominio acotado de CAA con 232 pictogramas admite pocas formulaciones válidas para la mayoría de las entradas, lo que acerca incluso a técnicas muy distintas hacia la misma solución.

Un hallazgo específico de la encuesta de naturalidad merece mención por su relevancia práctica. En el tipo comunicativo \textit{necesidades básicas}, el LLM no recibe ni una sola elección humana (0\% sobre 99 respuestas) frente al 42\% de los embeddings y el 39\% de la gramática, y un 18\% de los evaluadores marca ``Ninguna''. Para secuencias cortas y urgentes (``agua'', ``cansado'', ``baño''), los hablantes prefieren la formulación directa de los motores deterministas a la reformulación neuronal. Una arquitectura de producción puede enrutar según el tipo de secuencia, manteniendo al LLM como motor general pero conservando la gramática para enunciados donde la economía expresiva importa más que la elaboración sintáctica.

El desarrollo del proyecto ha consumido aproximadamente nueve meses, distribuidos entre el análisis del estado del arte (dos meses), la implementación de las tres técnicas base (tres meses), la hibridación y evaluación comparativa (dos meses), el desarrollo de la aplicación móvil (un mes) y la redacción de la memoria (un mes, solapado con las fases anteriores). Adicionalmente, el diseño, despliegue y recogida de datos de las dos encuestas con hablantes nativos han consumido alrededor de seis semanas efectivas entre la ejecución de la infraestructura (aplicación web, base de datos Supabase, cumplimiento RGPD) y el reclutamiento y seguimiento de participantes.

\section{Trabajo futuro}

La validación con usuarios reales de CAA es la línea de continuación más inmediata. Las dos encuestas ejecutadas en este TFG recogieron preferencias y referencias de hablantes nativos generales, pero no de personas con trastorno del espectro autista ni de profesionales del ámbito terapéutico o educativo. Un estudio con usuarios finales y con logopedas aportaría información sobre aceptación, usabilidad y adecuación al uso real que las métricas y las encuestas actuales no capturan. El hallazgo del 0\% de preferencia del LLM en \textit{necesidades básicas} sugiere además que el enrutamiento por tipo comunicativo debería validarse empíricamente con quienes usarían el sistema en primera persona.

La ampliación del vocabulario más allá de los 232 pictogramas actuales es el paso necesario para llevar el sistema a un entorno real. ARASAAC dispone de más de 40.000 pictogramas, y un vocabulario más amplio introduciría combinaciones que la gramática formal no puede anticipar con reglas manuales. En ese escenario, los embeddings semánticos podrían diferenciarse de la gramática al resolver predicciones que las reglas léxicas no cubren, y el LLM ganaría ventaja al manejar estructuras no previstas.

La evaluación con frases más largas y complejas complementaría la ampliación del vocabulario. El corpus actual contiene predominantemente frases de tres pictogramas; secuencias de seis o más pictogramas, con subordinación, coordinación múltiple o marcadores temporales, exigirían capacidades de generación que las reglas deterministas difícilmente pueden cubrir y que darían ventaja a los enfoques neuronales.

El ranking de Hybrid3 podría beneficiarse de un modelo de perplejidad más potente que \texttt{DeepESP/gpt2-spanish}. Un modelo con mayor capacidad de discriminación, por ejemplo un LLM moderno con instrucciones específicas de evaluación de naturalidad, podría seleccionar candidatas genuinamente mejores en los casos donde las frases divergen, en lugar de converger sistemáticamente con la gramática.

La integración completa del motor de generación en la aplicación móvil Dixi es el objetivo de producto natural del proyecto. Actualmente la app dispone de la interfaz de selección de pictogramas y del motor de síntesis de voz, pero la generación de frases se realiza offline con los motores de Python. Empaquetar la gramática formal como un servicio accesible desde la app, ya sea mediante una API o mediante un módulo nativo, completaría el ciclo de uso: el usuario selecciona pictogramas, el sistema genera una frase natural y el motor TTS la reproduce en voz alta.

Por último, la síntesis de voz personalizada constituye una línea de investigación abierta. Los motores TTS nativos de cada plataforma producen locuciones claras pero genéricas. La posibilidad de ajustar la voz a las características del usuario (edad, género, preferencia de tono) mejoraría la aceptación del sistema, dado que la voz es un componente central de la identidad comunicativa.
